\documentclass[11pt]{article}

\usepackage[margin=0.78in]{geometry}
\usepackage{amsmath,amssymb}
\usepackage{enumitem}

\setlength{\parindent}{0pt}
\setlength{\parskip}{4pt}

\begin{document}

{\large \textbf{COL 862/8595 }}

\vspace{3pt}

\textbf{Symbolic Execution}

How to know whether program satisfies requirements?

different possibilities:
model checking, testing, program analysis,
theorem proving, Hoare logic etc.

\vspace{7pt}

\textbf{testing}

Normal testing = concrete inputs.

For factorial, for example, can try

\[
n=0,\quad n=5,\quad n=-1,\quad n=20
\]

But only finite number of inputs get checked.

So passing all tests does not mean property is true
for every possible input.

\vspace{7pt}

\textbf{symbolic idea}

Instead of giving one fixed value to input,
give it a symbolic value.

\[
n \rightarrow \alpha_n
\]

Then execute the program using this symbolic value.

Main purpose = check whether a specified property
can ever be violated.

So it combines some ideas from testing + model checking.

\vspace{7pt}

For example:

\[
assert(result \geq 1 \lor n<0)
\]

Here $n$ is symbolic.

The execution engine keeps track of what must be
true for the current path.

\vspace{7pt}

\textbf{SEE}

Symbolic Execution Engine.

For every explored path/state, it keeps mainly:

\[
(\pi,\sigma,pc)
\]

where

\[
\pi = \text{path constraint}
\]

\[
\sigma = \text{symbolic store}
\]

\[
pc = \text{program counter}
\]

Path constraint is a first-order Boolean formula
which describes conditions needed to reach that path.

Symbolic store = values of variables may be expressions
instead of concrete numbers.

eg.

\[
x = \alpha_a + 2
\]

\vspace{7pt}

If assignment is

\[
x := e
\]

then symbolic store gets updated by substitution:

\[
\sigma[x\mapsto e_s]
\]

where $e_s$ is symbolic evaluation of $e$.

\vspace{7pt}

\textbf{branches}

Suppose current path constraint is $\pi$ and we reach

\[
if(e)\ then\ S_t\ else\ S_f
\]

two possible states are created.

true branch:

\[
\pi_t = \pi \land e_s
\]

false branch:

\[
\pi_f = \pi \land \neg e_s
\]

So one state can become two.

And then each of those can again branch.

This is where the number of paths starts getting big.

\vspace{7pt}

\textbf{solver part}

Need to know whether a path constraint is actually
possible.

Use SAT / SMT solvers.

\[
SAT \Rightarrow \text{path feasible}
\]

solver can also give concrete values for the symbols.

\[
UNSAT \Rightarrow \text{path impossible}
\]

\[
UNKNOWN/timeout \Rightarrow \text{couldn't determine}
\]

So symbolic execution can actually produce a concrete
test input corresponding to a path.

\vspace{7pt}

\textbf{constraints}

Modern engines use solvers such as Z3, CVC4, STP.

SAT asks basically:

\[
\text{is this formula satisfiable?}
\]

SMT extends this idea with theories.

For example arrays can be represented symbolically.

If

\[
B = Store(A,3,99)
\]

then $B$ is like $A$ except position $3$ gets value $99$.

For other positions:

\[
j\neq i \Rightarrow B[j]=A[j]
\]

\vspace{7pt}

\textbf{problem --- path explosion}

Every input-dependent branch can fork execution.

So number of states can become huge.

Main sources:

\[
\text{branches + loops + function calls}
\]

Loops are especially troublesome since iterations
can keep generating new paths.

Hence symbolic execution is powerful but cannot just
explore everything blindly.

\vspace{7pt}

\textbf{concolic}

In practice symbolic + concrete execution are often
mixed.

called

\[
\boxed{\text{concolic execution}}
\]

(dynamic symbolic execution is a common approach)

Concrete execution drives the symbolic exploration.

Maintain a concrete store too:

\[
\sigma_c : Var \rightarrow Val
\]

Suppose initially

\[
a=1,\qquad b=1
\]

Concrete execution follows one path.

The symbolic constraints for this path could be

\[
\pi=(\alpha_a\neq0)\land(\alpha_b\neq0)
\]

To explore another path, negate one branch condition
and ask the solver for a new input.

For example:

\[
(\alpha_a\neq0)\land(\alpha_b=0)
\]

gives one possible input

\[
a=1,\quad b=0
\]

which can drive execution through a different path.

\vspace{7pt}

Important point:

each concrete execution can reveal more unexplored
branches.

But SEE still has to decide which branch to explore next.

There can be a very large number of non-taken branches,
so heuristics matter.

\vspace{7pt}

\textbf{external code}

What if program calls some function whose implementation
is not available to SEE?

Then output of that function may determine a later branch.

SEE can sometimes reason about the path condition,
but generating an input that actually causes the external
function to return the required value is not guaranteed.

So external calls can cause problems in symbolic execution.

\vspace{7pt}

Another issue = path divergence.

Generated input was supposed to follow one symbolic path,
but when actually executed it follows another path.

Reason can be inaccurate modelling.

Examples mentioned include:

\[
\text{external calls, exceptions, casts, pointers}
\]

This can result in false negatives.

\vspace{7pt}

\textbf{choosing paths}

Since can't explore everything, choose paths intelligently.

Some methods:

\begin{itemize}[leftmargin=18pt,itemsep=1pt]
    \item DFS -- go deep first, then backtrack
    \item BFS -- expand paths level by level
    \item random path selection
\end{itemize}

KLEE can assign probabilities using path length
and branch arity.

Coverage based search can also give weights to states
and use them for selecting which state to explore.

\vspace{7pt}

\textbf{online / offline}

Online executor:

multiple paths explored in one run.

At an input-dependent branch,
state is cloned.

Advantage = don't repeat already executed instructions.

But memory usage can become large.

Offline executor:

one path at a time.

Less memory, but previous work may have to be repeated.

Checkpointing can reduce this repetition.

\vspace{7pt}

\textbf{memory}

Symbolic execution also has to model memory.

Memory can be thought of as mapping addresses to
concrete or symbolic expressions.

For a symbolic address, there can be multiple possible
locations.

eg. if

\[
a[i]
\]

and $i$ is symbolic, it may mean different array locations.

One option = fork the state for different possible
addresses.

So symbolic memory itself can create more states.

\vspace{7pt}

For arrays, theory of arrays can be used by the solver.

Full symbolic memory gives a more accurate model,
but can become intractable because of the explosion.

Another option is address concretization:

convert symbolic pointer/address to one particular
concrete address.

This reduces complexity but loses some precision.

\vspace{7pt}

\textbf{ITE}

Instead of always forking, possible values can also
be represented using an if-then-else expression.

\[
ITE(c,T,F)
\]

means

\[
\begin{cases}
T & c\text{ true}\\
F & c\text{ false}
\end{cases}
\]

Solvers can reason about these expressions without
necessarily creating a new state for every possibility.

\vspace{7pt}

\textbf{environment}

Programs interact with files, libraries, system calls etc.

SEE needs some model of these external effects.

For example KLEE has abstract models for external
functions and a basic symbolic file system.

Operations on symbolic files can again create branches.

\vspace{7pt}

\textbf{too many paths}

Ways to reduce this:

Prune paths which are impossible.

At a branch, check feasibility early.

If

\[
\pi=(x>0)\land(y=x+1)\land(z<0)\land(z>10)
\]

then the constraint

\[
z<0\land z>10
\]

is already impossible.

So paths containing the same contradiction can be
discarded.

This is where unsat cores are useful.

\vspace{7pt}

\textbf{summaries}

If same function or loop is executed many times,
doing symbolic exploration again and again is expensive.

Idea:

\[
\text{build summary} \rightarrow \text{save it}
\rightarrow \text{reuse}
\]

For every path, relate input conditions with output
conditions.

Then function summary can be a disjunction of the
possible path formulas.

Example:

\[
\phi_{isPos}
=
(x>0\land ret=1)
\lor
(x\leq0\land ret=0)
\]

This summary is independent of the calling context,
so it can be reused.

\vspace{7pt}

One issue though:

a concrete state can satisfy more than one precondition,
so using a summary does not always mean that the exact
same concrete path will be followed.

\vspace{7pt}

\textbf{LLM + symbolic execution}

LLMs can help in different parts of SE.

\begin{itemize}[leftmargin=18pt,itemsep=1pt]
    \item path prioritisation
    \item loop summarisation
    \item harness synthesis
    \item environment modelling
    \item approximate models for difficult computations
    \item explaining counterexamples / root cause
\end{itemize}

For path prioritisation, vulnerabilities can be used
to decide which paths deserve more attention.

For loops, LLMs may help produce a closed-form summary.

For external libraries / syscalls / callbacks,
LLMs can help construct approximate models.

\vspace{7pt}



\[
\text{symbolic inputs}
\rightarrow
\text{execute}
\rightarrow
\text{path constraints}
\rightarrow
\text{solver}
\]

then solver tells us whether the path is possible,
and if possible it can give an actual input.

Main difficulty is that the number of possible paths
and symbolic states can become very large.

\end{document}